% 负责人：A / hycx233 汇总，三人核对自己的交付与贡献。
\section{讨论与结论}
\subsection{已有结果支持的结论}
共享流程在六个真实样本上统一了坐标、深度归一化、距离背景、有效掩码及空间分组，
分类、发现、轨道与条件比较因此能够沿用同一套输入。所有量化结果均有保存的表格或
运行记录，代表图与全量生成结果分别管理。

分类最终测试使用 B 按原参数独立复跑的权重，在 104 个窗口上得到准确率 0.5096、
Macro-F1 0.3454；多数类基线分别为 0.7308、0.2815，逻辑回归为 0.8365、0.5644。
CNN 的宏平均 F1 高于多数类基线，但两项指标均低于逻辑回归，三个方法的 CHID F1 均为零。
这说明更复杂的模型在当前小样本和空间划分上没有必然优势，少数类和跨位置泛化仍是局限。
训练、测试预测、基线与 Grad-CAM 记录位于
\url{docs/results/classifier/test_evaluation/}；组长原权重的验证成绩单独保留，不与此次测试混写。

发现流程的 47 个高分窗口实际只对应 10 个独立位置，且重建误差与接触强度的
Pearson 相关为 0.828。非噪声簇中含有大量已知参考，未满足本组疑似新簇规则，
因此目前能够报告的是候选筛选、已知结构召回与候选排除，不能宣称发现了新结构。
聚类隐向量的空间分组交叉验证结果仅用于表征对照，不替代 CNN 的测试成绩。

WT 跨重复验证将 47 个重叠候选合并为 10 个独立位置，并按预先固定的规则各取一个
rep1 代表窗口。10 个代表窗口的共同有效上三角 O/E Pearson 相关为 0.6638--0.9158，
中位数 0.8572；4 个位置有无效区域，只能解释共同有效部分。8 个位置与已知注释重叠，
另 2 个不满足原新簇规则，因此没有可信新结构。这里的代表窗口结果不等同于 47 个窗口
逐一验证，也不证明结构类别或功能。

全基因组 465 张轨道图覆盖了最后不足 10 kb 的区间。已知结构的两种条件比较
共有 688 行可计算结果，分别有 217 和 182 条结构被标为一致增强；
同时保留减弱、参照带内与边界组合。CHIN\_126 和 OPCID\_53 在
$\Delta stpA$ 下分别展示了重复一致的增强与减弱，说明局部响应并非统一方向。

\subsection{局限与定稿依赖}
已知结构类别数为 68、250、26，类别分布不均，测试集中少数类结构也较少；
评价需要同时报告逐类样本量与 Macro-F1。同一基因组的空间划分有助于避免重叠输入
泄漏，但结果仍局限于当前数据与条件。自编码器的训练背景只有 17 个窗口，
其重建误差明显受接触强度影响，候选也不能只凭高分解释为新的形态类别。

条件比较使用每百万接触的相对量，增强不能推出细胞内绝对接触总量增加。
两个重复和固定参照带只支持描述性比较；双缺失的变化不能全部归于单个因子。
12 条结构的共同有效覆盖低于 80\%，对应 24 行结果仅反映共同有效区域，
故未用于代表案例。

候选跨重复验证已完成并按上述范围更新保留与排除结论；不据此扩大选做分析范围。
最终定稿仍需完成另一环境的代表流程复跑，并核对真实成员信息和贡献比例。

\section{成员分工与贡献}
下表记录计划分工，\textbf{不代表最终实际贡献比例}。
最终根据实际交付由三人共同确认百分比，总和为 100\%。

\begin{table}[htbp]
\centering
\begin{tabularx}{\linewidth}{l X c c}
\toprule
成员 & 主要分工 & 计划份额 & 实际贡献 \\
\midrule
A：hycx233 & 数据准备、共享预处理、差异定位、协助分类训练、整合与提交 & 约三分之一 & 待填写 \\
B：onlysonder & 分类与解释、重复验证、复跑与交付打包 & 约三分之一 & 待填写 \\
C：lanshuheng9-oss & 候选发现、聚类、多轨道图与案例 & 约三分之一 & 待填写 \\
\bottomrule
\end{tabularx}
\end{table}

三人分别撰写自己模块的报告章节、演示内容和讲稿；视频的实际录制安排由组长后续确定。
演示文稿统一使用 Beamer 编辑，默认使用 PDF 播放录制；PPTX 导出作为备用，最终提交讲解视频。
